\section{Decision-gated thesis programme}
\label{sec:programme}

\subsection{Provisional thesis claim}

The evidence supports the following provisional claim:

\begin{quote}
On a Xeon Max~9480 socket, the performance value of on-package HBM for
quantized LLM inference depends on workload phase, verified live-set
size and locality.  A controlled comparison can identify where HBM
changes TTFT and sustained decode, where its capacity becomes limiting,
and whether finer software placement is justified.
\end{quote}

This claim is narrower than ``HBM CPUs are better for inference'' and
does not assume that the runtime is the bottleneck.  It remains valuable
if HBM helps only decode, if the gain is modest, or if current software
fails to exploit the available bandwidth, provided the mechanism is
measured convincingly.

\subsection{Contribution ladder}

The programme is organised as a sequence of decisions.  Each stage
produces an independent artefact and determines whether the next stage
is warranted.

\begin{figure}[htbp]
\centering
\begin{tikzpicture}[
  >=Latex,
  stage/.style={
    draw=black!55,
    fill=neutralcol,
    rounded corners=2pt,
    text width=2.35cm,
    minimum height=1.45cm,
    align=center,
    font=\small
  },
  core/.style={stage, fill=accent!7, draw=accent},
  cond/.style={stage, fill=hbmcol!15, draw=hbmcol!80!black},
  note/.style={font=\scriptsize, align=center, text width=2.35cm}
]
  \node[stage] (s0) {Stage 0\\platform and harness};
  \node[core, right=4mm of s0] (s1) {Stage 1\\same-socket HBM/DDR};
  \node[core, right=4mm of s1] (s2) {Stage 2\\capacity and KV boundary};
  \node[stage, right=4mm of s2] (s3) {Stage 3\\topology policies};
  \node[cond, right=4mm of s3] (s4) {Stage 4\\selective tiering};

  \draw[->, thick] (s0) -- (s1);
  \draw[->, thick] (s1) -- (s2);
  \draw[->, thick] (s2) -- (s3);
  \draw[->, thick] (s3) -- (s4);

  \node[note, below=5mm of s1] {core thesis after a sound\\phase-resolved result};
  \node[note, below=5mm of s4] {only if benefit, capacity\\pressure and missing control coexist};
\end{tikzpicture}
\caption{Contribution ladder.  Later stages are options unlocked by
evidence, not commitments made in advance.}
\label{fig:programme}
\end{figure}

\subsection{Stage 0: platform and measurement qualification}

The first deliverable is a dated platform record and a benchmark
harness.  It should capture:

\begin{itemize}
  \item allocated node identity, BIOS/firmware, kernel, scheduler
        request, CPU frequency policy and memory/clustering mode;
  \item the NUMA distance matrix, CPU lists, installed and usable HBM
        and DDR capacity, transparent huge-page state and automatic
        NUMA-balancing state;
  \item compiler, \lcpp{} commit, build flags, loaded modules, model
        hash, tokenizer, quantization and command line;
  \item local-HBM, local-DDR and deliberately remote STREAM/latency
        measurements under the intended binding; and
  \item a machine-readable manifest linked to each raw result.
\end{itemize}

The stage passes when repeated microbenchmarks distinguish local HBM,
local DDR and remote memory in the expected direction, and page
residency confirms the requested policy.  If the system cannot expose
Flat mode or distinguish tiers, the research question must be revised
with the supervisors before inference runs begin.

\subsection{Stage 1: primary same-socket experiment}

Stage 1 answers \researchq{1} with an unmodified, pinned runtime.  The smallest
credible matrix varies:

\begin{itemize}
  \item memory policy: local HBM and paired local DDR;
  \item phase: fixed-prompt prefill and steady-state autoregressive
        decode;
  \item at least two model/quantization points that fit comfortably in
        usable local HBM; and
  \item a small set of batch or sequence counts that changes arithmetic
        intensity without exhausting capacity.
\end{itemize}

The primary outcomes are TTFT, ITL and phase throughput.  Mechanism
measurements include resident pages by NUMA node, peak resident memory,
memory bandwidth, instructions or cycles, CPU frequency and remote
traffic.  Repetitions are interleaved or randomised across memory
policies to reduce drift.  The analysis reports effect sizes and
confidence intervals, not only the best run.

Stage 1 is already a defensible empirical contribution if placement is
verified and the result explains both positive and null effects.  It
also supplies the data needed to decide whether the rest of the
programme is worthwhile.

\subsection{Stage 2: capacity and context}

Stage 2 answers \researchq{2}.  Models, quantizations, context lengths and
concurrent sequences are chosen to approach the measured usable HBM
capacity from below and to cross it in controlled increments.  Before
performance is interpreted, the experiment records:

\[
  M_{\mathrm{live}} =
  M_{\mathrm{weights}} + M_{\mathrm{packed}} + M_{\mathrm{KV}}
  + M_{\mathrm{graph/scratch}} + M_{\mathrm{other}} .
\]

The terms should be measured where possible.  Formula-based KV
estimates are retained as checks, not substituted for observed peak
allocation.  A capacity curve plots the performance effect against
verified live-set fraction \(M_{\mathrm{live}}/M_{\mathrm{HBM,usable}}\),
with failures and fallback behaviour shown rather than discarded.

If context changes alter the token sequence or output-quality setting,
those changes are recorded.  Performance comparisons must use the same
model representation and decoding policy.

\subsection{Stage 3: topology policies}

Stage 3 starts only after the one-socket mechanism is stable.  Its
purpose is not to maximise core count; it is to test whether explicit
locality changes the scaling conclusion.  The preferred comparison is
one socket versus two socket-local replicas.  A spanning-process
configuration is a negative control.  SNC4-specific work is included
only if the mode is available and its CPU/memory pairings can be
verified.

The stage passes academically even when replicas do not improve
latency: aggregate throughput, fairness, remote traffic and
per-request performance together establish the trade-off.  A complex
tensor-parallel CPU implementation is outside this stage unless an
existing supported path can be configured without code changes.

\subsection{Stage 4: conditional runtime contribution}

Stage 4 answers \researchq{4} only when the preceding results show:

\begin{enumerate}
  \item a material HBM advantage in a named phase;
  \item a repeatable capacity limit;
  \item a policy reason to keep different data classes in different
        tiers; and
  \item no existing, auditable mechanism that expresses that policy.
\end{enumerate}

The first implementation target is a distinct, selectable CPU HBM
buffer type using the runtime's existing abstractions.  The second,
only if required, is a KV-cache buffer selector.  Each change needs unit
or integration tests for allocation, clear failure on insufficient HBM,
page-residency validation and an ablation against the best unmodified
policy.  Upstream compatibility and a small patch are preferable to a
fork-specific scheduler.

No performance claim follows merely from making a control available.
The systems contribution is the combination of a justified abstraction,
correct placement and a measured benefit or well-explained negative
result.

\subsection{Deliverables and stopping rules}

\begin{table}[htbp]
\centering
\small
\caption{Deliverables and explicit stopping rules.}
\label{tab:deliverables}
\begin{tabularx}{\textwidth}{@{}P{2.25cm}YY@{}}
\toprule
\textbf{Stage} & \textbf{Deliverable} & \textbf{Stop or redirect when} \\
\midrule
0: qualify
& Platform manifest, placement checks, local/remote memory baselines
& HBM/DDR cannot be controlled or observed on allocated nodes. \\
\addlinespace
1: isolate
& Phase-resolved HBM/DDR dataset and causal analysis
& No stable inference baseline exists; fix reproducibility before
  expanding the matrix. \\
\addlinespace
2: cross capacity
& Live-set and performance boundary with failures retained
& The accessible models never approach usable HBM; choose a larger
  model/context or report the verified range. \\
\addlinespace
3: topology
& Named placement-policy comparison with remote-traffic evidence
& Required modes/counters are unavailable or the result would not
  answer a distinct question. \\
\addlinespace
4: implement
& Minimal selective-tier control, tests and ablation
& Existing controls suffice, no capacity pressure exists, or selective
  placement has no credible mechanism. \\
\bottomrule
\end{tabularx}
\end{table}

The thesis is complete after Stages 0--2 if they are executed rigorously.
Stage 3 strengthens the locality contribution.  Stage 4 changes the
thesis from primarily empirical to empirical-plus-systems, but it is not
required for a coherent result.

\subsection{Risks to manage early}

\begin{itemize}
  \item \textbf{Operational access:} memory modes, counters or node
        exclusivity may be restricted.  Confirm them before selecting
        the final experiment matrix.
  \item \textbf{Page-cache contamination:} file-backed weights can
        retain earlier placement.  Develop a non-destructive,
        administrator-approved cold/warm protocol and verify pages.
  \item \textbf{Thermal and frequency drift:} interleave treatments,
        log frequency and temperature where available, and avoid
        comparing runs from different node states without blocking.
  \item \textbf{Runtime movement:} pin the commit and dependencies for
        the main study.  Treat a later runtime as a separate replication,
        not an in-place update.
  \item \textbf{Matrix growth:} use pilot data to select informative
        points.  More configurations do not compensate for weak
        placement evidence or too few repetitions.
\end{itemize}
